\section*{写作说明与协作分工}
\addcontentsline{toc}{section}{写作说明与协作分工}
课程要求报告包含需求分析、概要设计、详细设计和系统测试。数据库部分要有教材式 ER，讲清规范化与开发便利性的取舍，并审查大模型辅助设计的结果\cite{course}。
各人编辑自己负责的文件，写作提示用于查漏。可以调整小节顺序，选自己熟悉的例子，把图放在它要解释的段落旁。

\begin{table}[htbp]\centering\small
\begin{tabularx}{\linewidth}{lYY}
\toprule 成员角色 & 写在哪里 & 图放在哪里\\\midrule
A 数据库 & 数据库设计、成长与协作、模型辅助审查 & ER 与字段示例放建模，宽表分解放规范化，查询对照放查询设计\\
B 界面 & 需求分析、系统实现与交互 & 场景图放需求，真截图按操作顺序贯穿系统交互\\
C 答疑与检索 & 概要设计、答疑、历史与资料、记忆、摘要与总结 & 总架构图放概要设计，时序图和引用图放对应模块\\
D 测试与集成 & 判题、系统测试、部署与运行维护 & 判题图放练习，性能图放实验结果，部署图放迁移与启动\\\bottomrule
\end{tabularx}
\caption{各人负责的章节与配图}
\end{table}

各人把做过的事、作图源稿和运行证据一起交给 D。写正文时可以从一个操作讲起，把改动、理由和结果写具体。图注交代读者看到什么，真实截图注明实际操作与运行模式。

各人的章文件见 \path{docs/report/report-plan.md}。素材和此前的配图建议见 \path{docs/report/report-plan-detailed.md}。新图稿放 \path{assets/}。合稿后把本说明和各处提示换成正式正文。
